\section{Biểu diễn logic tự động và chuyển đổi sang CNF}

\subsection{Biến đổi các manh mối quan hệ hai ngôi}

Mọi lời khai đang mở đều được biên dịch tự động sang CNF thay vì viết sẵn công thức theo màn. Với \texttt{SAME}, quan hệ tương đương $C_A\leftrightarrow C_B$ được tách thành hai chiều kéo theo rồi áp dụng $p\to q\equiv \neg p\lor q$:

\[
\begin{aligned}
\operatorname{SAME}(A,B)&\equiv C_A\leftrightarrow C_B\\
&\equiv(C_A\to C_B)\land(C_B\to C_A)\\
&\equiv(\neg C_A\lor C_B)\land(C_A\lor\neg C_B).
\end{aligned}
\]

Với \texttt{DIFFERENT}, điều kiện hai người trái trạng thái tương đương có ít nhất một người là tội phạm và không thể cả hai cùng là tội phạm:

\[
\begin{aligned}
\operatorname{DIFFERENT}(A,B)&\equiv C_A\oplus C_B\\
&\equiv (C_A\lor C_B)\land\neg(C_A\land C_B)\\
&\equiv (C_A\lor C_B)\land(\neg C_A\lor\neg C_B).
\end{aligned}
\]

\subsection{Ví dụ mã hóa đếm tổ hợp trực tiếp trên vùng nhỏ}

Xét vùng ba biến $\{A,B,C\}$ với yêu cầu đúng $k=2$ tội phạm. Để có ít nhất 2 tội phạm thì mọi cặp 2 ô đều phải chứa ít nhất một biến đúng; để có nhiều nhất 2 thì cả 3 ô không thể cùng đúng:

\[
\begin{aligned}
\operatorname{EXACTLY}(2,\{A,B,C\})\equiv{}&
(A\lor B)\land(A\lor C)\land(B\lor C)\\
&{}\land(\neg A\lor\neg B\lor\neg C).
\end{aligned}
\]

Tổng quát trên vùng $n$ ô, mã hóa tổ hợp trực tiếp sinh $\binom{n}{n-k+1}=\binom{n}{k-1}$ mệnh đề dương và $\binom{n}{k+1}$ mệnh đề âm, cộng lại là $\binom{n}{k-1}+\binom{n}{k+1}$ mệnh đề (bằng 0 ngoài miền chỉ số). Với $n=25,\ k=12$, riêng một lời khai đã cần $9\,657\,700$ mệnh đề nên không phù hợp cho bàn lớn.

\clearpage
\begingroup\setstretch{1.25}

\subsection{Quy hoạch động đếm ngưỡng unary và cổng Tseitin}

Để giữ kích thước CNF bậc đa thức $O(n^2)$ và tái sử dụng cho phép so sánh vùng, \texttt{Encoder} lập bảng quy hoạch động các ngưỡng unary $T(i,k)$ — mệnh đề “ít nhất $k$ trong $i$ ô đầu của vùng là tội phạm”, với $T(i,0)=\mathrm{True}$ và $T(i,k)=\mathrm{False}$ khi $k>i$. Khi thêm ô thứ $i$ ($C_i$), ngưỡng $k$ đạt được nếu $i-1$ ô trước đã đạt $k$, hoặc đã đạt $k-1$ và $C_i$ đúng:

\[
T(i,k)\leftrightarrow T(i-1,k)\lor\bigl(C_i\land T(i-1,k-1)\bigr).
\]

Mỗi ô $(i,k)$ tạo tối đa một cổng \texttt{AND} và một cổng \texttt{OR} hai đầu vào. Biến phụ được cấp từ $N^2+1$ trong một \texttt{Encoder} chung cho cả $KB_t$; hàm \texttt{gate()} rút gọn hằng, khử literal lặp, triệt tiêu cặp đối dấu và tra bảng băm để dùng lại cổng trùng nhau.

Theo phép biến đổi Tseitin \cite{tseitin1983complexity}, đầu ra phụ được ràng buộc hai chiều để an toàn cả dưới phép phủ định. Với $z\leftrightarrow(a\land b)$ và $w\leftrightarrow(a\lor b)$, ta có:

\[
\begin{aligned}
&(\neg z\lor a)\land(\neg z\lor b)\land(z\lor\neg a\lor\neg b),\\
&(w\lor\neg a)\land(w\lor\neg b)\land(\neg w\lor a\lor b).
\end{aligned}
\]

<<<<<<< HEAD
Đặt $T_R(k)=T(|R|,k)$. Ba ràng buộc “ít nhất”, “nhiều nhất” và “đúng bằng” lần lượt nhận literal hoặc biểu thức $T_R(k)$, $\neg T_R(k+1)$ và $T_R(k)\land\neg T_R(k+1)$. Trong cài đặt thực tế, cả các manh mối hai ngôi như \texttt{SAME} và \texttt{DIFFERENT} cũng được biểu diễn qua \texttt{Encoder.gate} để kết hợp trực tiếp bên trong manh mối điều kiện. Bộ mã hóa buộc đầu ra của mỗi manh mối công khai bằng một mệnh đề đơn; đầu ra hằng sai sinh mệnh đề rỗng, còn đầu ra hằng đúng không cần thêm mệnh đề.
=======
Đặt $T_R(k)=T(|R|,k)$. Ba mẫu \texttt{AT\_LEAST}, \texttt{AT\_MOST} và \texttt{EXACTLY} lần lượt trả về $T_R(k)$, $\neg T_R(k+1)$ và $T_R(k)\land\neg T_R(k+1)$. Đầu ra mỗi lời khai công khai được ép bằng một mệnh đề đơn (nếu hằng sai thì thêm mệnh đề rỗng, hằng đúng thì bỏ qua).
>>>>>>> e95a2af30fd2e650472e82590bb55816b145db80

\subsection{Mã hóa các manh mối mở rộng và kiểm chứng độc lập}

Trong \texttt{COMPARISON}, số tội phạm của $R_1$ lớn hơn hoặc bằng $R_2$ khi mọi ngưỡng $k\ge 1$ đạt ở $R_2$ đều đạt ở $R_1$:

\[
\bigwedge_{k=1}^{|R_2|}\bigl(\neg T_{R_2}(k)\lor T_{R_1}(k)\bigr).
\]

Hoán vị hai vùng cho $\le$, lấy hội hai chiều cho $=$, và phủ định để được $\ne$, $<$ hoặc $>$. Với \texttt{CONDITIONAL}, bộ mã hóa đệ quy lấy hai đầu ra $r_{\mathrm{if}},r_{\mathrm{then}}$ rồi ghép cổng $\neg r_{\mathrm{if}}\lor r_{\mathrm{then}}$. Độc lập với CNF, mỗi lớp đều có hàm \texttt{evaluate()} tính trực tiếp trên phép gán đầy đủ để kiểm thử chéo.

\par\endgroup
